\chapter{The expression language}\label{app:expr}\label{ch:expr}

\begin{goals}
\textbf{In\index{expression language} this appendix}: the whole expression language in one place — how names, numbers and text
are written, every operator and function, what a blank does, and a cookbook of expressions you can
paste and adapt.
\end{goals}

micro has one expression language and uses it everywhere a condition or a formula is needed:

\begin{itemize}
\item the filter box and its drawer (\cref{ch:filter});
\item ƒ calculated columns (\cref{ch:filter});
\item the density and weight slots and the per-series filters of the analysis windows
  (\cref{ch:analysis});
\item the hole filter, over collar columns (\cref{ch:drillholes});
\item rule-based styling (\cref{ch:color});
\item in a routine: the \code{filter} of an \code{each} block, the computed values written
  \code{\$\{…\}}, and the checks (\cref{ch:routines,ch:checks}).
\end{itemize}

An expression reads the columns of \emph{one record} — one block, one interval, one row — and
gives back a value: a number, a piece of text, true or false, or a blank. Used as a filter, the
record is kept when the value is true.

\begin{why}
The language is \term[expression language!total]{total}: an expression can compare and compute, but it cannot loop, call
anything outside itself, or stop a scan with an error. Dividing by zero, taking the logarithm of a
negative number or doing arithmetic on text all give a blank. That is what makes it safe to open a
project, a recipe or a routine that somebody else wrote: an expression in it can be wrong, but it
cannot do anything other than produce values.
\end{why}

\section{Writing an expression}\label{sec:expr-syntax}

\begin{lstlisting}[language=expr]
FE > 55 and LITO = "HEMATITE"
FE between 40 and 55 or LITO in ("CANGA", "ITABIRITE")
FE / 100 * SG * 500        # tonnes of iron in a 10 × 10 × 5 m block
\end{lstlisting}

A bare word is a column. A quoted word is text. So \code{FE > SG} compares two columns, and
\code{LITO = "HEMATITE"} tests a column against a value. Forgetting the quotes
(\code{LITO = HEMATITE}) asks for a column named \col{HEMATITE}, and micro says that no such column
exists.

\subsection{Names}\label{sec:expr-names}

\begin{tabularx}{\linewidth}{@{}p{42mm}L@{}}
\toprule
you write & it means \\
\midrule
\code{FE}, \code{fe}, \code{Fe\_pct} & A column, by name. Letters, digits and underscores, not
  starting with a digit. Case doesn't matter: \code{fe} finds \col{FE}. \\
\code{Fe.pct}, \code{g.cutoff} & A dotted name is still one name. The dot is part of the name;
  nothing is being looked up inside anything. An assay column called \col{Fe.pct} and a routine's
  \code{g.cutoff} are written the same way. \\
\code{\textasciigrave Assay Au ppm\textasciigrave} & Backticks, for a name with spaces, hyphens or
  other symbols, or a name that is also a keyword (\code{\textasciigrave in\textasciigrave}). Inside
  backticks, write \code{\textbackslash\textasciigrave} for a backtick and
  \code{\textbackslash\textbackslash} for a backslash. \\
\code{\$\{fe-pct\}} & The same thing in the brace form: letters, digits, \code{\_}, \code{.} and
  \code{-} between the braces. A routine uses this form for its computed values
  (\cref{sec:expr-routines}). \\
\bottomrule
\end{tabularx}

\begin{caution}
A hyphen is always a minus sign. \code{OK-INDIC} is the column \col{OK} minus the column
\col{INDIC}. For a column that is really named \col{OK-INDIC}, write
\code{\textasciigrave OK-INDIC\textasciigrave}; micro suggests exactly that when it sees the bare
form and no such columns.
\end{caution}

The words \code{and}, \code{or}, \code{not}, \code{between}, \code{contains}, \code{in},
\code{like}, \code{matches}, \code{is}, \code{blank}, \code{filled}, \code{true} and \code{false}
are reserved. Keywords and function names are not case-sensitive (\code{AND}, \code{If(…)}). A
function name is a function only when a parenthesis follows it, so a column that happens to be
called \col{round} or \col{min} still works as a bare name.

In the color list and the record panel a calculated column appears as \ui{ƒ NAME}, and a stored
one — materialized, joined, painted, a fraction inside a solid — with a ✎ or a \%. Those marks
are labels, not part of the name: in an expression, write the name alone.

\subsection{Numbers, text and comments}\label{sec:expr-literals}

\begin{tabularx}{\linewidth}{@{}p{42mm}L@{}}
\toprule
you write & it means \\
\midrule
\code{42} \; \code{3.14} \; \code{.5} \; \code{1.5e-3} & Numbers. Always with a decimal point, never
  a comma. \\
\code{"HEMATITE"} & Text, in double quotes. Single quotes are accepted too. \\
\code{true} \; \code{false} & The two truth values. \\
\code{blank} & The blank value itself, as in \code{if(FE = -99, blank, FE)}. \\
\code{\# to the end of the line} & A comment. An expression may run over several lines, so a long
  filter can be laid out and annotated. \\
\bottomrule
\end{tabularx}

\begin{note}
\textbf{Decimal commas.} A file whose numbers are written \code{1,5} (the Brazilian and European
convention) is detected when it opens, and its columns are read as numbers everywhere
(\cref{ch:opening}). The numbers you \emph{type} in an expression still take a point:
\code{FE > 55.5}, never \code{FE > 55,5}. A comma in an expression separates the arguments of a
function and the members of an \code{in (…)} list.
\end{note}

\section{Operators}\label{sec:expr-operators}

\begin{tabularx}{\linewidth}{@{}p{26mm}p{42mm}L@{}}
\toprule
kind & operators & notes \\
\midrule
Arithmetic & \code{+} \; \code{-} \; \code{*} \; \code{/} & Division by zero gives a blank. \\
Power & \code{\textasciicircum{}} (or \code{**}) & \code{FE\textasciicircum{}2}. \\
Comparison & \code{=} \; \code{!=} \; \code{<} \; \code{<=} \; \code{>} \; \code{>=} & Numbers
  compare as numbers; text compares as text, and exactly: \code{"Hematite"} is not
  \code{"HEMATITE"}. \\
Range & \code{x between a and b} & Both ends included. \\
Membership & \code{x in (a, b, c)} & True when \code{x} is one of the listed values. \\
Text search & \code{x contains "txt"} & The text appears somewhere in the value. \\
 & \code{x like "DH\%\_1"} & SQL wildcards over the whole value: \code{\%} is any run of
  characters, \code{\_} is exactly one. \\
 & \code{x matches "\textasciicircum{}DH[0-9]+"} & A regular expression. \\
Blank tests & \code{x is blank} \; \code{x is filled} & Also \code{is not blank},
  \code{is not filled}. \\
Negated forms & \code{not in (…)} \; \code{not contains} \; \code{not like} & \\
Logic & \code{and} \; \code{or} \; \code{not} & With parentheses for grouping. \\
\bottomrule
\end{tabularx}

Text comparisons, \code{in}, \code{contains}, \code{like} and \code{matches} are case-sensitive. To
compare without regard to case, bring the column to one case first:
\code{upper(LITO) = "HEMATITE"}.

A few spellings from other languages are accepted, so old habits still work: \code{==} for
\code{=}, \code{<>} for \code{!=}, \code{\&\&} and \code{||} for \code{and} and \code{or},
\code{\textasciitilde{}} for \code{contains}. The forms in the table are the ones micro highlights and suggests.

\begin{caution}
A \code{matches} pattern is checked before it is used. A pattern that could take unbounded time on
a crafted value — a repeated group that itself holds a repeat or an alternative, such as
\code{(a+)+} or \code{(a|b)*}, or a back-reference — is refused, as is a pattern that isn't valid.
Rewrite it without the nesting: \code{[ab]*} does what \code{(a|b)*} does.
\end{caution}

\subsection{Precedence}\label{sec:expr-precedence}

From the loosest binding to the tightest. Operators on one line bind equally and group from the
left, except the power, which groups from the right.

\begin{tabularx}{\linewidth}{@{}p{8mm}p{52mm}L@{}}
\toprule
 & operators & example \\
\midrule
1 & \code{or} & \\
2 & \code{and} & \code{a or b and c} is \code{a or (b and c)} \\
3 & \code{not} & \code{not a and b} is \code{(not a) and b} \\
4 & comparisons, \code{between}, \code{in}, \code{contains}, \code{like}, \code{matches},
  \code{is blank} & \code{FE + 1 > SG * 2} compares the two sums \\
5 & \code{+} \; \code{-} & \\
6 & \code{*} \; \code{/} & \code{1 + 2 * 3} is 7 \\
7 & the sign: \code{-x} & \\
8 & \code{\textasciicircum{}} & \code{-2\textasciicircum{}2} is −4;
  \code{2\textasciicircum{}-3} is 0.125; \code{2\textasciicircum{}3\textasciicircum{}2} is 512 \\
9 & names, numbers, text, function calls, parentheses & \\
\bottomrule
\end{tabularx}

\begin{tip}
\code{not} takes one comparison. To negate more, use parentheses:
\code{not (FE > 55 and LITO = "HEMATITE")}. When in doubt about \code{and} and \code{or}, add the
parentheses; the filter drawer shows them as indented groups, which makes the structure visible.
\end{tip}

\section{Functions}\label{sec:expr-functions}

Every function gives a blank when an argument it needs is blank or of the wrong kind, unless the
table says otherwise. Square brackets mark an optional argument.

\subsection{Choosing}

\begin{tabularx}{\linewidth}{@{}p{46mm}L@{}}
\toprule
function & gives \\
\midrule
\code{if(cond, a, b)} & \code{a} when the condition is true, otherwise \code{b}. Nest it for a
  ladder of cases. Only the chosen branch is computed. \\
\code{coalesce(a, b, …)} & The first argument that isn't blank. \\
\code{ifnum(x, d)} & \code{x} as a number when it is one, otherwise \code{d}. \\
\code{nullif(x, v)} & Blank when \code{x} equals \code{v}, otherwise \code{x}: the way to scrub a
  sentinel such as −99. \\
\bottomrule
\end{tabularx}

\subsection{Rounding and binning}

\begin{tabularx}{\linewidth}{@{}p{46mm}L@{}}
\toprule
function & gives \\
\midrule
\code{round(x[, n])} & \code{x} rounded to \code{n} decimals (none by default). \\
\code{int(x)} & \code{x} with its decimals cut off, toward zero: \code{int(-2.7)} is −2. \\
\code{floor(x)} \; \code{ceil(x)} & The whole number below, or above: \code{floor(-2.7)} is −3. \\
\code{abs(x)} & The absolute value. \\
\code{mod(x, y)} & The remainder, always with the sign of \code{y}, as in a spreadsheet:
  \code{mod(-7, 3)} is 2. Blank when \code{y} is 0. \\
\code{clamp(x, lo, hi)} & \code{x}, held between \code{lo} and \code{hi}. \\
\code{bin(x, width[, origin])} & The lower edge of the bin of that width that holds \code{x}, with
  bins starting at \code{origin} (0 by default): \code{bin(1047, 10)} is 1040, and
  \code{bin(1047, 10, 5)} is 1045. Blank when the width is zero or negative. \\
\bottomrule
\end{tabularx}

\subsection{Math}

\begin{tabularx}{\linewidth}{@{}p{46mm}L@{}}
\toprule
function & gives \\
\midrule
\code{sqrt(x)} & The square root. Blank for a negative number. \\
\code{pow(x, y)} & \code{x} to the power \code{y}; the same as
  \code{x\textasciicircum{}y}. \\
\code{exp(x)} & $e^x$. \\
\code{log(x)} & The natural logarithm. Blank for zero or less. \\
\code{log10(x)} & The base-10 logarithm, the usual transform for skewed grades. \\
\code{min(a, b, …)} \; \code{max(a, b, …)} & The smallest, or largest, of the numeric arguments.
  Blanks are skipped; blank only when every argument is blank. \\
\code{lerp(a, b, t)} & The point a fraction \code{t} of the way from \code{a} to \code{b}:
  $a + (b - a)\,t$. \code{t} is not held to 0–1: \code{lerp(2, 4, 0.25)} is 2.5, and
  \code{lerp(2, 4, 2)} is 6. \\
\code{remap(x, a1, b1, a2, b2)} & \code{x} rescaled so that \code{a1}…\code{b1} lands on
  \code{a2}…\code{b2}: \code{remap(40, 20, 60, 0, 1)} is 0.5. Not held to the range (wrap it
  in \code{clamp} for that); blank when \code{a1} = \code{b1}. \\
\code{interp(x, x1, y1, x2, y2, …)} & The straight-line curve through the points
  (\code{x1}, \code{y1}), (\code{x2}, \code{y2}), …, read at \code{x}, and flat beyond the
  first and last: \code{interp(55, 20, 2.7, 45, 3.4, 65, 4.3)} is 3.85. Needs at least two
  points. Points written as numbers must have increasing \code{x1}, \code{x2}, …, or the
  expression is refused; points taken from columns are put in order for you. \\
\code{smoothstep(a, b, x)} & A smooth 0-to-1 ramp as \code{x} goes from \code{a} to
  \code{b}, 0 before and 1 after: \code{smoothstep(0, 10, 2.5)} is 0.15625. \\
\bottomrule
\end{tabularx}

\subsection{Tests}

\begin{tabularx}{\linewidth}{@{}p{46mm}L@{}}
\toprule
function & true when \\
\midrule
\code{isblank(x)} \; \code{isfilled(x)} & \code{x} is blank; \code{x} is not blank. The same as
  \code{x is blank} and \code{x is filled}. \\
\code{isnum(x)} & \code{x} is a number, or text that reads as one. \\
\code{isnan(x)} & \code{x} is filled but isn't a number: a stray \code{n/a} in a grade column. \\
\bottomrule
\end{tabularx}

These four always answer true or false, never blank.

\subsection{Text}

\begin{tabularx}{\linewidth}{@{}p{46mm}L@{}}
\toprule
function & gives \\
\midrule
\code{upper(s)} \; \code{lower(s)} & The text in capitals; in lower case. \\
\code{trim(s)} & The text without spaces at either end. \\
\code{len(s)} & The number of characters. \\
\code{left(s, n)} \; \code{right(s, n)} & The first, or last, \code{n} characters. \\
\code{substr(s, start[, len])} & The text from position \code{start}, counting from 1, for
  \code{len} characters (to the end by default). \\
\code{replace(s, find, with)} & The text with every \code{find} replaced. \code{find} is taken
  literally, not as a pattern. A blank \code{with} deletes. \\
\code{concat(a, b, …)} & The arguments joined. Blanks are skipped; numbers are written as text:
  \code{concat("DH", 7)} is \code{"DH7"}. \\
\bottomrule
\end{tabularx}

Text that comes out empty is a blank: \code{trim("  ")} is blank, and \code{len} of a blank is
blank, not 0.

\subsection{Dates}

\begin{tabularx}{\linewidth}{@{}p{46mm}L@{}}
\toprule
function & gives \\
\midrule
\code{year(d)} \; \code{month(d)} \; \code{day(d)} & The part of a date written
  \code{YYYY-MM-DD} (anything after the day is ignored). Blank for any other spelling. \\
\bottomrule
\end{tabularx}

\section{Types}\label{sec:expr-types}

A value is a number, text, true or false, or blank. You rarely have to think about which, because
the language converts where the meaning is clear:

\begin{itemize}
\item Text that reads as a number is a number in arithmetic and in a comparison with a number. A
  text column holding \code{"007"} equals \code{7}.
\item When either side of a comparison isn't a number, both are compared as text. Ordering text
  (\code{<}, \code{>}) is by character order, so \code{"DH10" < "DH9"} is true.
\item Arithmetic on text that isn't a number gives a blank, never an error and never 0.
\item \code{in (…)} compares the written form, so \code{BENCH in (1040, 1050)} works on a number
  column and on a text column alike.
\item true and false are not numbers: \code{(FE > 55) + 1} is blank. To count, write
  \code{if(FE > 55, 1, 0)}.
\end{itemize}

The Calculations window shows a calculated column's type as a choice (\ui{auto}, \ui{number},
\ui{text}); see \cref{ch:filter}. A source column's type is shown,
and can be changed, in \menu{Properties>columns} (\cref{ch:columns}).

\section{Blanks}\label{sec:expr-blanks}

A \term[blank (missing value)]{blank} is a missing value: an empty cell, a null in a Parquet file, a lab code read as no
value, a number that isn't one. micro never turns a blank into zero by itself, because a missing
grade that silently becomes 0 drags every mean and every tonnage with it.

The rules:

\begin{itemize}
\item \textbf{Blanks propagate.} Arithmetic and functions with a blank input give a blank:
  \code{FE * SG} is blank where either is. The exceptions are the functions made for blanks —
  \code{coalesce}, \code{ifnum}, the tests — and \code{min}, \code{max} and \code{concat}, which
  skip them.
\item \textbf{An ordering with a blank is not true.} Where \col{FE} is blank, \code{FE > 55} and
  \code{FE <= 55} are both not true, so a filter on either leaves the record out.
\item \textbf{Blank equals blank.} \code{FE = blank} is true where \col{FE} is blank, and
  \code{A = B} is true where both are. There is no SQL-style trap where a comparison with a
  missing value matches nothing; still, \code{FE is blank} is the clearer spelling.
\item \textbf{\code{not} is a complement.} A filter and its negation always split the table in two
  with nothing left over.
\item \textbf{\code{if} takes the else branch} when its condition is blank or false.
\end{itemize}

\begin{tabularx}{\linewidth}{@{}p{50mm}p{20mm}L@{}}
\toprule
where \col{FE} is blank & result & \\
\midrule
\code{FE > 55} & not true & a blank never satisfies an ordering \\
\code{FE <= 55} & not true & \\
\code{not (FE > 55)} & true & the complement takes it \\
\code{FE between 40 and 55} & not true & \\
\code{FE = blank} \; \code{FE is blank} & true & \\
\code{FE != 55} & true & a blank is not 55 \\
\code{FE + 1} & blank & \\
\code{if(FE > 55, 1, 0)} & 0 & the else branch \\
\code{coalesce(FE, 0)} & 0 & you asked for the zero \\
\bottomrule
\end{tabularx}

\begin{caution}
\code{not (FE > 55)} and \code{FE <= 55} are different filters: the first keeps the records with a
blank \col{FE}, the second doesn't. The same goes for \code{FE != 55}, which keeps blanks. When you
mean “assayed, and at most 55”, say so: \code{FE is filled and FE <= 55}.
\end{caution}

\begin{caution}
\code{if(FE > 55, "ore", "waste")} calls an unassayed record “waste”, because a blank condition
takes the else branch. Test for the blank first:
\code{if(FE is blank, blank, if(FE > 55, "ore", "waste"))}.
\end{caution}

Sentinels are not blanks until you say so. A file that writes −99 for “not assayed” has a column of
numbers, and its mean includes them. Either tell micro once, for a column of a CSV or drillhole table
(\menu{Properties>columns}, right-click the column, \ui{Detection limits \& blanks};
\cref{ch:columns}), or
write \code{nullif(FE, -99)} where you use it.

\section{In routines and checks}\label{sec:expr-routines}

A routine uses the same language in three places (\cref{ch:routines,ch:checks}).

\begin{lstlisting}[language=yaml]
- each:
    from: "params.xlsx#Benches"
  as: "b"
  filter: "active = 1 and z0 >= g.floor"
  steps:
    - recipe: "deposit grade-tonnage"
      with:
        cutoffs:
          - '${if(b.lito = "HEMATITE", 55, 45)}'
      as: "gt"
- check: "abs(val.bias_pct) < 5"
\end{lstlisting}

\begin{description}
\item[The \code{filter} of an \code{each} block] decides which rows of the parameter table run.
  The row's columns are plain names; scalars from \code{params} are visible under their namespace
  (\code{g.floor}).
\item[\code{\$\{…\}}] computes a value from the names in scope. On its own it gives a typed value (a
  number stays a number); inside longer text it is written into the text. A bare \code{\$name}
  substitutes a value as it stands, with no computing.
\item[A check] reads the values of earlier steps, each under the step's \code{as} name
  (\code{val.bias\_pct}), and the baseline run's under \code{prev.} (\code{prev.val.bias\_pct}). A
  row check (\code{every:}) reads the columns of a step's table by name, and \code{prev.tonnes} for
  the same row of the baseline.
\end{description}

This is where dotted names matter: \code{b.bench}, \code{g.cutoff} and \code{val.bias\_pct} are
each one name. A mistyped one is reported with the nearest real name, never read as empty.

\begin{note}
A check that reads a blank is \emph{not evaluated}. It is not a pass. Since blank equals blank in
this language, a check comparing two missing values would otherwise succeed for the wrong reason;
micro reports which names were blank and, at level \code{stop}, stops the run.
\end{note}

\section{A cookbook}\label{sec:expr-cookbook}

The examples use the demo project's columns where they can: the block model \file{deposit.csv}
(\col{XC}, \col{YC}, \col{ZC}, \col{FE}, \col{SG}, \col{LITO}) and the drillholes (\col{HOLEID},
\col{COMPANY}, \col{EOH} on the collars; \col{HOLEID}, \col{FROM}, \col{TO}, \col{FE}, \col{SG}
and \col{LITO} on the assays). Columns such as \col{AU},
\col{MN} or \col{DOMAIN} stand for your own.

\subsection{Filters}

\begin{tabularx}{\linewidth}{@{}p{50mm}L@{}}
\toprule
to keep & write \\
\midrule
High-grade hematite & \code{FE > 60 and LITO = "HEMATITE"} \\
A grade band & \code{FE between 45 and 55} \\
Two lithologies & \code{LITO in ("HEMATITE", "ITABIRITE")} \\
Everything but waste and canga & \code{LITO not in ("WASTE", "CANGA")} \\
One bench (10 m, toe at 600) & \code{ZC between 600 and 610} \\
A box in plan & \code{XC between 612500 and 613500 and YC between 7765300 and 7765700} \\
Blocks with no estimate & \code{FE is blank} \\
Estimated blocks above cutoff & \code{FE is filled and FE >= 45} \\
Ore by either of two rules & \code{(FE > 58) or (FE > 50 and LITO = "HEMATITE")} \\
Lithology codes starting with ITA & \code{LITO like "ITA\%"} \\
Holes DH001 to DH009 & \code{HOLEID matches "\textasciicircum{}DH00[1-9]\$"} \\
One company's holes (hole filter) & \code{COMPANY = "PIONEER"} \\
Short holes (hole filter) & \code{EOH < 200} \\
Values sitting exactly on a cap & \code{FE = 68} \\
Junk in a grade column & \code{isnan(AU)} \\
\bottomrule
\end{tabularx}

\subsection{Derived columns}

\begin{tabularx}{\linewidth}{@{}p{50mm}L@{}}
\toprule
to get & write \\
\midrule
Sentinel −99 as blank & \code{nullif(AU, -99)} \\
A capped grade & \code{clamp(AU, 0, 30)} \\
The first filled of two estimates & \code{coalesce(AU\_OK, AU\_ID2)} \\
An equivalent grade & \code{FE + 0.4 * MN} \\
Tonnes per block (10 × 10 × 5 m) & \code{SG * 500} \\
Contained metal per block, t & \code{FE / 100 * SG * 500} \\
Density by lithology & \code{if(LITO = "HEMATITE", 3.9, if(LITO = "ITABIRITE", 3.4, 2.8))} \\
Density from grade, along a curve & \code{interp(FE, 20, 2.7, 45, 3.4, 65, 4.3)} \\
A grade rescaled to 0–1 between 20 and 65 & \code{clamp(remap(FE, 20, 65, 0, 1), 0, 1)} \\
A default density where none is given & \code{coalesce(SG, 2.7)} \\
An ore / waste flag that keeps blanks blank &
  \code{if(FE is blank, blank, if(FE >= 45, "ORE", "WASTE"))} \\
Grade classes as a code & \code{if(FE > 58, 2, if(FE > 45, 1, 0))} \\
Bench toe from elevation (10 m benches) & \code{bin(ZC, 10)} \\
Bench number counted from 400 m & \code{floor((ZC - 400) / 10)} \\
A 50 m panel index along east & \code{bin(XC, 50)} \\
Log grade, for a skewed color scale & \code{log10(clamp(AU, 0.01, 100))} \\
Relative difference of two estimates, \% & \code{(AU\_OK - AU\_NN) / AU\_NN * 100} \\
Absolute difference & \code{abs(AU\_OK - AU\_NN)} \\
Grade rounded for a report & \code{round(FE, 1)} \\
Distance from a point in plan & \code{sqrt((XC - 613000)\textasciicircum{}2 + (YC -
  7765500)\textasciicircum{}2)} \\
A ratio that is blank, not infinite, on zero & \code{SIO2 / AL2O3} \\
\bottomrule
\end{tabularx}

The last line makes a point: nothing needs guarding. Where \col{AL2O3} is 0 or blank the
ratio is blank, and the statistics leave that record out.

\subsection{Drillholes and text}

\begin{tabularx}{\linewidth}{@{}p{50mm}L@{}}
\toprule
to get & write \\
\midrule
Interval length & \code{TO - FROM} \\
Clean hole ids for a join & \code{upper(trim(HOLEID))} \\
Hole number from DH007 & \code{ifnum(substr(HOLEID, 3), blank)} \\
A sample key & \code{concat(HOLEID, "-", FROM, "-", TO)} \\
Ids with a hyphen removed & \code{replace(HOLEID, "-", "")} \\
The drilling year & \code{year(DATE)} \\
Intervals longer than 2 m (a filter) & \code{TO - FROM > 2} \\
A column with a space in its name & \code{\textasciigrave Fe pct\textasciigrave{} > 55} \\
\bottomrule
\end{tabularx}

\subsection{Routines and checks}

\begin{tabularx}{\linewidth}{@{}p{50mm}L@{}}
\toprule
to say & write \\
\midrule
A cutoff 10\% above the sheet's & \code{\$\{g.cutoff * 1.1\}} \\
A cutoff that depends on the row & \code{\$\{if(b.lito = "HEMATITE", 55, 45)\}} \\
Only the active rows of a table & \code{filter: "active = 1"} \\
Global bias within 5\% & \code{check: "abs(val.bias\_pct) < 5"} \\
Every row has tonnes & under \code{check:}, \code{every: "tonnes > 0"} and \code{in: "gt"} \\
No row moved more than 10\% from the baseline &
  under \code{check:}, \code{every: "abs(tonnes / prev.tonnes - 1) < 0.1"}, \code{in: "gt"} and
  \code{against: "prev"} \\
\bottomrule
\end{tabularx}

Each expression in the first three tables works anywhere an expression is accepted: the filter box,
a ƒ column, a density or weight slot, a per-series filter, the hole filter.
